\documentclass[11pt,a4paper]{article}
\usepackage[T1]{fontenc}
\usepackage[utf8]{inputenc}
\usepackage[ngerman,english]{babel}
\usepackage{lmodern}
\usepackage{microtype}
\usepackage{geometry}
\geometry{a4paper, margin=2.5cm}
\usepackage{graphicx}
\usepackage{xcolor}
\usepackage{hyperref}
\hypersetup{
  colorlinks=true,
  linkcolor=blue!70!black,
  urlcolor=blue!60!black,
  citecolor=green!60!black,
  pdftitle={EtherCAT Driver ROS2 -- Technical Report},
  pdfauthor={Jaron Martinez}
}
\usepackage{listings}
\usepackage{booktabs}
\usepackage{longtable}
\usepackage{array}
\usepackage{tabularx}
\usepackage{float}
\usepackage{caption}
\usepackage{subcaption}
\usepackage{enumitem}
\usepackage{amsmath}
\usepackage{amssymb}
\usepackage{tcolorbox}
\tcbuselibrary{skins,breakable}
\usepackage{mdframed}
\usepackage{tikz}
\usetikzlibrary{arrows.meta,positioning,shapes,fit,backgrounds,calc}
\usepackage{pgfplots}
\pgfplotsset{compat=1.18}
\usepackage{fancyhdr}
\usepackage{titlesec}
\usepackage{parskip}

% ----------- Colors ------------------------------------------------------------
\definecolor{codegreen}{rgb}{0,0.6,0}
\definecolor{codegray}{rgb}{0.5,0.5,0.5}
\definecolor{codepurple}{rgb}{0.58,0,0.82}
\definecolor{backcolour}{rgb}{0.97,0.97,0.97}
\definecolor{warnorange}{RGB}{230,120,0}
\definecolor{dangerred}{RGB}{200,30,30}
\definecolor{safegreen}{RGB}{30,150,30}
\definecolor{infoblue}{RGB}{30,100,180}

% ----------- Listings ----------------------------------------------------------
\lstdefinestyle{cppstyle}{
  backgroundcolor=\color{backcolour},
  commentstyle=\color{codegreen}\itshape,
  keywordstyle=\color{blue!70!black}\bfseries,
  numberstyle=\tiny\color{codegray},
  stringstyle=\color{codepurple},
  basicstyle=\ttfamily\small,
  breakatwhitespace=false,
  breaklines=true,
  captionpos=b,
  keepspaces=true,
  numbers=left,
  numbersep=5pt,
  showspaces=false,
  showstringspaces=false,
  showtabs=false,
  tabsize=2,
  language=C++,
  morekeywords={uint8_t,uint16_t,uint32_t,int8_t,int16_t,int32_t,size_t,constexpr,override,nullptr,std}
}
\lstset{style=cppstyle}

% ----------- Custom boxes ------------------------------------------------------
\newtcolorbox{bugbox}[2][]{
  colback=red!5!white,colframe=dangerred,
  fonttitle=\bfseries,title=Bug #2,#1
}
\newtcolorbox{fixbox}[1][]{
  colback=green!5!white,colframe=safegreen,
  fonttitle=\bfseries,title=Fix,#1
}
\newtcolorbox{infobox}[1][]{
  colback=blue!5!white,colframe=infoblue,
  fonttitle=\bfseries,#1
}
\newtcolorbox{warnbox}[1][]{
  colback=orange!8!white,colframe=warnorange,
  fonttitle=\bfseries,#1
}

% ----------- Header / Footer ---------------------------------------------------
\pagestyle{fancy}
\fancyhf{}
\rhead{\small\texttt{ethercat\_driver\_ros2} -- Technical Report}
\lhead{\small Jaron Martinez, 2026}
\cfoot{\thepage}
\renewcommand{\headrulewidth}{0.4pt}

% ----------- Title formatting --------------------------------------------------
\titleformat{\section}{\Large\bfseries\color{blue!70!black}}{}{0em}{\thesection\quad}
\titleformat{\subsection}{\large\bfseries\color{blue!50!black}}{}{0em}{\thesubsection\quad}
\titleformat{\subsubsection}{\normalsize\bfseries\color{blue!40!black}}{}{0em}{\thesubsubsection\quad}

% ==============================================================================
\begin{document}

\begin{titlepage}
  \centering
  \vspace*{2cm}
  {\Huge\bfseries EtherCAT Driver ROS\,2\\[0.4em]
   \Large Technical Report: Branch Evolution, Safety Architecture,\\
   and Bug-Fix Analysis\par}
  \vspace{1.5cm}
  {\large\itshape LeRobot Teststand Project\par}
  \vspace{0.5cm}
  {\large Jaron Martinez\par}
  \vspace{0.5cm}
  {\large\texttt{jaron.martinez@gmail.com}\par}
  \vspace{1cm}
  {\large 28 June 2026\par}
  \vspace{2cm}
  \hrule
  \vspace{1em}
  \begin{abstract}
    \noindent
    This report documents the full technical evolution of the
    \texttt{ethercat\_driver\_ros2} fork developed for Synapticon Circulo
    servo drives (CiA402) controlled via an IgH EtherCAT Master and a
    Beckhoff CU1228 coupler in a LeRobot robotic-arm teststand.
    Starting from the upstream \texttt{main} branch of ICube-Robotics,
    two major feature branches were created:
    \texttt{saftey\_features\_clean} (immediate safety stop on link or
    working-counter fault) and \texttt{deterministic-barrier-startup}
    (deterministic Distributed-Clock synchronisation and a group
    barrier that eliminates serial per-joint initialisation delays).
    The report covers the motivation for each change, the precise code
    modifications made, and a rigorous analysis of eleven bugs
    discovered and fixed in the barrier-startup branch, including three
    IEEE\,754 NaN-safety issues.  A commissioning and test plan
    concludes the document.
  \end{abstract}
  \hrule
  \vspace{1em}
  \tableofcontents
\end{titlepage}

% ==============================================================================
\newpage
\section{Introduction and Motivation}
\label{sec:intro}

The \texttt{ethercat\_driver\_ros2} package provides a \emph{ros2\_control}
\texttt{SystemInterface} that connects ROS\,2 hardware controllers to EtherCAT
fieldbus devices via the IgH EtherCAT Master for Linux.  The upstream project
(ICube-Robotics, University of Strasbourg) is a general-purpose driver
suitable for any EtherCAT slave.  The fork documented here specialises it for
a six-degree-of-freedom robotic arm equipped with Synapticon Circulo drives,
which implement the \emph{CANopen over EtherCAT} (CoE) CiA402 motion-control
profile.

\subsection{Hardware Setup}

\begin{itemize}[nosep]
  \item \textbf{PC}: Real-time Linux (PREEMPT\_RT), IgH EtherCAT Master 1.5
  \item \textbf{EtherCAT coupler}: Beckhoff CU1228 (bus switch)
  \item \textbf{Drives}: 6$\times$ Synapticon Circulo (vendor ID 0x000022D2,
        CiA402 Cyclic Synchronous Position mode, 1\,kHz PDO cycle)
  \item \textbf{ROS\,2}: Humble, \texttt{ros2\_control}, \texttt{pluginlib}
\end{itemize}

\subsection{Branch Overview}

\begin{center}
\begin{tikzpicture}[
  node distance=2.2cm,
  commit/.style={rectangle,rounded corners,draw,minimum width=4cm,
                 minimum height=0.7cm,font=\small\ttfamily,fill=blue!8},
  arrow/.style={-Stealth,thick,blue!60}
]
  \node[commit] (main)   {\textbf{main} (upstream ICube)};
  \node[commit,below=of main] (safety)
        {\textbf{saftey\_features\_clean}};
  \node[commit,below=of safety] (barrier)
        {\textbf{deterministic-barrier-startup}};

  \draw[arrow] (main)   -- node[right,font=\small]{Apr 2026 -- safety stop} (safety);
  \draw[arrow] (safety) -- node[right,font=\small]{Jun 2026 -- DC fix + barrier + 11 bugs} (barrier);
\end{tikzpicture}
\end{center}

% ==============================================================================
\section{Starting Point: Upstream \texttt{main}}
\label{sec:upstream}

The upstream \texttt{main} branch provided:
\begin{enumerate}[nosep]
  \item A YAML-configurable plugin loader (\texttt{pluginlib}) for any
        \texttt{EcSlave} subclass.
  \item A generic PDO channel manager (\texttt{EcPdoChannelManager}) with
        read/write, masks, bit operations, and a NaN firewall
        (\texttt{ec\_write} skips the write when \texttt{default\_value} is
        NaN).
  \item A CiA402 drive plugin (\texttt{EcCiA402Drive}) with automatic state
        transitions, mode-of-operation control, and fault-reset edge detection.
  \item The IgH \texttt{EcMaster} wrapper (PDO domain management, SDO
        download, DC configuration).
\end{enumerate}

\subsection{Known Issues in Upstream}
\label{sec:upstream-issues}

Even in the upstream release the following weaknesses existed:

\begin{warnbox}[title=Upstream limitations]
  \begin{itemize}[nosep]
    \item No link-down detection in the real-time loop: a cable pull caused
          the driver to continue silently sending stale commands.
    \item No WC-failure counter: a single working-counter mismatch generated a
          \texttt{WARN} log but no controlled stop.
    \item \texttt{initialized\_} was set to \texttt{is\_operational\_}
          (EtherCAT bus-OP flag) rather than to the \emph{drive} state
          \texttt{OPERATION\_ENABLED}, so the driver could declare success before
          any drive was actually ready to accept motion commands.
    \item All drive state interfaces were initialised to NaN and were never
          written to a defined value until the first PDO cycle returned data,
          creating a window where \texttt{last\_position\_} was NaN.
    \item Per-slave \texttt{clock\_gettime()} in \texttt{addSlave()} produced
          random, non-deterministic SYNC0 phases, causing 1--5\,s per-joint
          DC-alignment delays during startup.
  \end{itemize}
\end{warnbox}

% ==============================================================================
\section{Branch 1: \texttt{saftey\_features\_clean}}
\label{sec:safety-branch}

This branch was created in April 2026 to address the most critical safety gaps
identified during initial teststand commissioning.

\subsection{Commit \texttt{b21391f} -- Immediate ERROR on Link-Down / WKC Fault (Apr 2026)}

\subsubsection{Motivation}

During hardware testing, pulling the EtherCAT cable while the driver was
active caused the control loop to continue running with stale, non-updated
PDO data.  The drives continued executing the last received target position,
which could cause uncontrolled motion.

\subsubsection{Changes}

The \texttt{read()} callback was extended with two sequential safety checks:

\begin{lstlisting}[caption={Link-down detection in \texttt{read()}}]
// SAFETY: Check link status -- immediate ERROR on link down
if (!master_->isMasterLinkUp()) {
  RCLCPP_ERROR(..., "SAFETY: EtherCAT link DOWN -- stopping ...");
  activated_ = false;
  consecutive_wc_failures_ = 0;
  master_->deactivate();
  return hardware_interface::return_type::ERROR;
}

// SAFETY: Check domain WKC -- immediate ERROR when not COMPLETE
unsigned int wc_state = master_->getDomainWcState(0);
if (wc_state != EC_WC_COMPLETE) {
  consecutive_wc_failures_++;
  if (consecutive_wc_failures_ >= kMaxConsecutiveWcFailures_) {
    RCLCPP_ERROR(..., "SAFETY: Domain WC not COMPLETE for %d "
                      "consecutive cycles -- returning ERROR!",
                 consecutive_wc_failures_, wc_state);
    return hardware_interface::return_type::ERROR;
  }
}
\end{lstlisting}

The \texttt{isMasterLinkUp()} method wraps \texttt{ecrt\_master\_state()} and
checks the \texttt{link\_up} field.  Returning \texttt{ERROR} from
\texttt{read()} causes \texttt{ros2\_control} to call
\texttt{on\_error()}, which calls \texttt{master\_->deactivate()} and then
requests a clean process shutdown via \texttt{rclcpp::shutdown()}.

\subsubsection{Initial WC Threshold: \texttt{kMaxConsecutiveWcFailures = 1}}

The initial value of 1 was overly aggressive: a single late EtherCAT frame
during startup or a momentary CPU overload could trigger a safety stop.

\subsection{Commit \texttt{70f0595} -- Tighten Safety Stop and Clean Shutdown (Apr 2026)}

Added \texttt{on\_error()} to the hardware interface lifecycle to guarantee a
controlled motor stop in all error paths:

\begin{lstlisting}[caption={\texttt{on\_error()} safe shutdown}]
CallbackReturn EthercatDriver::on_error(
  const rclcpp_lifecycle::State &)
{
  if (activated_) {
    activated_ = false;
    master_->deactivate();  // slaves -> INIT, brakes engage
  }
  // Delayed shutdown to let controller_manager finish
  if (rclcpp::ok() && !shutdown_requested_.exchange(true)) {
    std::thread([]() {
      std::this_thread::sleep_for(std::chrono::milliseconds(200));
      if (rclcpp::ok()) { rclcpp::shutdown(); }
    }).detach();
  }
  return CallbackReturn::SUCCESS;
}
\end{lstlisting}

\texttt{ecrt\_master\_deactivate()} followed by
\texttt{ecrt\_release\_master()} puts all slaves back to the \texttt{INIT}
state, which engages any hardware brakes configured on the drives.  The
subsequent \texttt{ecrt\_release\_master()} also prevents the IgH kernel
idle-scanning thread from re-establishing slave states after a cable
reconnect without explicit re-activation.

\subsection{Commit \texttt{32a2cc6} -- Hold Latched Position When Controller Inactive (Apr 2026)}

\subsubsection{Problem -- Floating-Point Accumulation Drift at 1\,kHz}

\paragraph{Observed symptom.}
When the trajectory controller is inactive (command interface set to NaN by
ROS\,2's hardware-interface layer), a joint that should be stationary slowly
drifts away from its resting position.  The drift is not caused by external
forces or a faulty encoder -- it is a purely software artefact that appears
only in standby mode and worsens over time.

\paragraph{Root cause.}
In Cyclic Synchronous Position (CSP) mode the drive expects a target position
command every EtherCAT cycle.  When no trajectory controller is active the
driver must write a \emph{hold position} to keep the joint stationary.  The
upstream implementation recomputed this hold value from scratch on \emph{every
single 1\,ms control cycle}:

\begin{lstlisting}[caption={Upstream hold-position code executed at 1\,kHz (buggy)}]
// Inside processData(), every 1 ms:
channel.default_value =
    channel.factor * last_position_ + channel.offset;
\end{lstlisting}

\noindent
\texttt{last\_position\_} is the actual position read back from the drive's
TPDO in the same cycle.  This forms a closed feedback loop:

\begin{enumerate}
  \item Cycle $n$: send $d_n = f \cdot p_n + b$ as the commanded position.
  \item Drive moves toward $d_n$; encoder reads $p_{n+1}$ (actual feedback).
  \item Cycle $n+1$: compute $d_{n+1} = f \cdot p_{n+1} + b$.
\end{enumerate}

\noindent
This loop is a problem for two independent reasons:

\begin{description}
  \item[IEEE\,754 rounding per multiply.]
    The factor $f$ is loaded from the YAML configuration as a double-precision
    float.  Typical unit-conversion factors (e.g.\ encoder counts to radians:
    $f = 2\pi / 1\,048\,576 \approx 5.99 \times 10^{-6}$, or its inverse for
    the command channel) are \emph{not} exactly representable in binary
    floating point.  Each multiplication $f \cdot p_n$ therefore introduces a
    rounding error of up to $\tfrac{1}{2}\,\mathrm{ULP}$.  For a 20-bit
    encoder with values around $10^6$, ULP$(10^6 \cdot 2\pi) \approx 10^{-9}$\,rad.
    Individually negligible, at 1\,kHz this sums to $\approx 10^{-6}$\,rad$/$s.
    Over minutes the error can accumulate to multiple encoder counts.

  \item[Integer quantisation round-trip.]
    The drive's servo loop operates on fixed-point integers (32-bit position
    counts from a 20-bit encoder).  The TPDO delivers the actual position as
    an integer-valued double ($p_n \in \mathbb{Z}$).  After the float multiply
    and round-trip through the RPDO the drive rounds the received command back
    to an integer count.  If $\lfloor f \cdot p_n + b \rceil \neq p_n$, the
    next TPDO sample $p_{n+1}$ can differ from $p_n$ by $\pm 1$ count,
    which is then fed back into the multiply loop.  This quantisation walk
    causes slow, monotone position drift -- typically 1--5 encoder
    counts per minute, corresponding to
    $\approx 0.001$--$0.02$\,degrees for a 20-bit encoder.
\end{description}

\paragraph{Why this is hard to catch in testing.}
The drift is small enough that a short activation/deactivation cycle does not
reveal it.  It only becomes observable after the driver has been running in
standby for 30--60 seconds, manifesting as an unexpected joint motion that
looks like a control or gravity-compensation bug rather than a driver artefact.
A \texttt{rqt\_plot} of the joint position command vs.\ state interface while
the controller is inactive will show the command slowly diverging.

\subsubsection{Solution -- One-Shot Position Latch with NaN Guard}

A \emph{position-hold latch} was introduced: when the controller command
interface is NaN (controller inactive), the drive's \texttt{default\_value}
for the RPDO position is set \emph{once} to the current actual position and
held constant until the controller becomes active again.  The
\texttt{position\_hold\_initialized\_} flag ensures the computation runs
exactly once per controller-deactivation event, completely eliminating the
feedback loop.

A subsequent fix (branch \texttt{deterministic-barrier-startup}, this session)
added a \texttt{std::isfinite()} guard before the latch commit.  Without this
guard, if the TPDO position had not yet been received when the latch triggered
(e.g.\ during the first few EtherCAT cycles after activation),
\texttt{last\_position\_} would be \texttt{NaN} and the drive would receive a
NaN command -- which the lower-level \texttt{ec\_write()} firewall in
\texttt{ec\_pdo\_channel\_manager.cpp} silently suppresses, leaving the domain
at its previous (potentially zero) value.

\begin{lstlisting}[caption={Position-hold latch in \texttt{processData()}}]
bool controller_active = false;
if (channel.has_command_interface_name(0) &&
    channel.is_command_interface_defined()) {
  const size_t cmd_idx = channel.command_interface_index(0);
  if (cmd_idx < command_interface_ptr_->size()) {
    controller_active =
      !std::isnan(command_interface_ptr_->at(cmd_idx));
  }
}
if (controller_active) {
  if (std::isfinite(last_position_)) {
    channel.default_value =
      channel.factor * last_position_ + channel.offset;
  }
  position_hold_initialized_ = false;
} else {
  // Latch once, then hold
  if (!position_hold_initialized_ && std::isfinite(last_position_)) {
    channel.default_value =
      channel.factor * last_position_ + channel.offset;
    position_hold_initialized_ = true;
  }
}
\end{lstlisting}

\subsection{Commit \texttt{4f1e593} -- Watchdog Tuning and Monotonic App Time (Apr 2026)}

The EtherCAT SM watchdog was configured explicitly for approximately 100\,ms
timeout (divider 2498, intervals 1000 at 25\,ns base period):

\begin{lstlisting}
ecrt_slave_config_watchdog(slave_info.config, 2498, 1000);
\end{lstlisting}

A fresh \texttt{CLOCK\_MONOTONIC} timestamp is passed to
\texttt{ecrt\_master\_application\_time()} on every \texttt{update()} call,
ensuring the DC reference clock stays synchronised during long-running
operation rather than drifting from a single reference set at activation.

\subsection{Commit \texttt{aedb118} -- Raise \texttt{kMaxConsecutiveWcFailures} to 5 (May 2026)}

Experience with the 1-cycle threshold showed spurious safety stops during
rapid controller switches and brief CPU scheduling jitter.  The threshold was
raised to 5 consecutive WC failures, allowing transient single-cycle glitches
to recover without a controlled stop.

% ==============================================================================
\section{Branch 2: \texttt{deterministic-barrier-startup}}
\label{sec:barrier-branch}

This branch, developed in June 2026, addresses the two remaining major
problems: non-deterministic DC startup delays and the absence of a
group-synchronised power-up gate.

\subsection{Commit \texttt{0ff09ed} -- Deterministic SYNC0 Shift (Jun 2026)}
\label{sec:sync0-fix}

\subsubsection{Root Cause of 20-Second Startup Times}

The upstream \texttt{addSlave()} function was called once per slave.  Inside
it, a \texttt{clock\_gettime(CLOCK\_MONOTONIC, \&t)} was issued for each
slave, and the SYNC0 shift was computed as:

\begin{lstlisting}[language=C,caption={Old (broken) per-slave SYNC0 shift}]
// OLD: issued per slave -- different timestamps, different phases
clock_gettime(CLOCK_MONOTONIC, &t);
uint32_t sync0_shift = interval_ - (t.tv_nsec % interval_);
ecrt_slave_config_dc(config, 0x0100, interval_, sync0_shift, 0, 0);
\end{lstlisting}

Because \texttt{addSlave()} was called sequentially for each of the six
drives, each received a \emph{different} effective SYNC0 phase.  During the
SAFEOP$\to$OP transition, the IgH master waited for each drive's individual
SYNC0 alignment before releasing it to OP, leading to roughly 1--5\,s per
joint and a total startup time of 15--20\,seconds.

\begin{center}
\begin{tikzpicture}
  \begin{axis}[
    xlabel={Joint index},
    ylabel={Startup time (s)},
    xtick={0,1,2,3,4,5},
    ymin=0, ymax=22,
    width=9cm, height=5cm,
    bar width=0.5cm,
    ybar,
    nodes near coords,
    every node near coord/.append style={font=\footnotesize},
    title={Simulated worst-case startup (random SYNC0 phase)},
    ymajorgrids=true,
    grid style={dotted}
  ]
    \addplot coordinates {(0,1.2)(1,3.4)(2,2.8)(3,5.1)(4,4.0)(5,1.9)};
    \addplot[fill=safegreen!50,draw=safegreen]
            coordinates {(0,0.09)(1,0.09)(2,0.09)(3,0.09)(4,0.09)(5,0.09)};
    \legend{Random phase (before fix), Fixed phase (after fix)}
  \end{axis}
\end{tikzpicture}
\end{center}

\subsubsection{Fix}

A single, fixed SYNC0 shift equal to half the cycle interval (cycle midpoint)
is applied \emph{identically} to every slave.  The application time is set
\emph{once} in \texttt{activate()} immediately before
\texttt{ecrt\_master\_activate()}:

\begin{lstlisting}[caption={Fixed SYNC0 shift in \texttt{addSlave()}}]
if (slave->assign_activate_dc_sync()) {
  dc_used_ = true;
  const uint32_t sync0_shift = interval_ / 2;  // same for all slaves
  ecrt_slave_config_dc(
    slave_info.config,
    slave->assign_activate_dc_sync(),
    interval_,
    sync0_shift,
    0, 0);
}
\end{lstlisting}

With all six drives sharing the same SYNC0 phase, the IgH master aligns the
entire bus simultaneously, reducing total startup time from $\sim$20\,s to
$\sim$0.5\,s (dominated by CiA402 state-machine transitions).

\subsection{Commit \texttt{a8268c7} -- Strict Init Readiness Gate (Jun 2026)}

The activation loop was redesigned with a strict gate that requires all
modules to be \emph{simultaneously} ready before declaring success.  The gate
uses three independent criteria:

\begin{enumerate}[nosep]
  \item \textbf{Domain WC state} = \texttt{EC\_WC\_COMPLETE} (all slaves OP
        and responding)
  \item \textbf{\texttt{readyForCommands()}} = \texttt{true}
        (CiA402 state machine in \texttt{OPERATION\_ENABLED})
  \item \textbf{\texttt{hasValidPosition()}} = \texttt{true}
        (at least one finite \texttt{TPDO\_POSITION} received)
\end{enumerate}

All three conditions must hold for \texttt{init\_stable\_cycles\_} (default
10) consecutive cycles.  If they do not hold within \texttt{init\_timeout\_}
seconds (default 30\,s), activation aborts and reports which module is stuck
along with its status string (CiA402 state + position validity + bus-OP flag).

\subsection{Commit \texttt{902be31} -- Align Phase-0 Barrier Timeout and Debounce (Jun 2026)}

Fine-tuned the \texttt{phase\_stable\_cycles\_} default from 20 to 10 cycles
(matching the WC check frequency) and logged the total expected wait at the
start of \texttt{on\_activate()}:

\begin{lstlisting}[caption={Total timeout log at activation start}]
RCLCPP_INFO(...,
  "Starting ...please wait... "
  "(max %.0f s barrier + %.0f s init = %.0f s total)",
  phase_timeout_, init_timeout_,
  phase_timeout_ + init_timeout_);
\end{lstlisting}

\subsection{Commit \texttt{3917d66} -- WIP: Fix 11 Bugs (Jun 2026)}
\label{sec:bugfixes}

A thorough review of the barrier-startup implementation against the original
\texttt{main} branch revealed eleven bugs.  Each is described with its
severity, impact, and fix below.

% ==============================================================================
\section{Bug Analysis and Fixes}
\label{sec:bugs}

\begin{table}[H]
  \centering
  \caption{Summary of all 11 fixed bugs}
  \label{tab:bugs}
  \small
  \begin{tabularx}{\textwidth}{@{}c l X l@{}}
    \toprule
    \# & Location & Description & Severity \\
    \midrule
    1 & \texttt{generic\_ec\_cia402\_drive.cpp} &
        \texttt{initialized\_} semantics: was set to EtherCAT bus-OP, not
        CiA402 OPERATION\_ENABLED & High \\
    2 & \texttt{ethercat\_driver.cpp} /
        \texttt{generic\_ec\_cia402\_drive.hpp} &
        Shadowed function name: static helper had same name as virtual
        method added later & High \\
    3 & \texttt{ec\_master.cpp} &
        No null-pointer check before \texttt{memcpy} in
        \texttt{transferAll()} & Medium \\
    4 & \texttt{ethercat\_driver.hpp} &
        \texttt{phase\_stable\_cycles\_} default 20 vs.\ 10 intended & Low \\
    5 & \texttt{ethercat\_driver.cpp} &
        Missing total-timeout log at \texttt{on\_activate()} entry & Low \\
    6 & \texttt{ethercat\_driver.cpp} &
        SDO failure logged at \texttt{INFO} instead of \texttt{ERROR} & Medium \\
    7 & \texttt{ethercat\_driver.cpp} &
        \texttt{std::stod} used for integer slave position parameter & Low \\
    8 & \texttt{ethercat\_driver.cpp} &
        Double push of transfer slaves into \texttt{ec\_transfer\_slaves\_} & Medium \\
    9 & \texttt{ethercat\_driver.hpp/.cpp} &
        \texttt{kMaxConsecutiveWcFailures} was \texttt{constexpr} (not
        configurable at runtime) & Medium \\
    NaN-1 & \texttt{generic\_ec\_cia402\_drive.cpp} &
        Fault-reset latch stuck on NaN command (IEEE\,754 trap) & High \\
    NaN-2 & \texttt{generic\_ec\_cia402\_drive.cpp} &
        Position-hold latch entered when \texttt{last\_position\_} is NaN & Medium \\
    NaN-3 & \texttt{generic\_ec\_cia402\_drive.cpp} &
        Controller-active path wrote NaN to \texttt{default\_value} & Medium \\
    \bottomrule
  \end{tabularx}
\end{table}

\subsection{Bug~1 -- \texttt{initialized\_} Semantics}
\label{sec:bug1}

\textbf{Location:} \texttt{generic\_ec\_cia402\_drive.cpp:64},
\texttt{updateState()}

\textbf{Impact:} The activation loop in \texttt{on\_activate()} checked
\texttt{module->initialized()} and \texttt{module->readyForCommands()}.
Because \texttt{initialized\_} was set to \texttt{is\_operational\_} (the
EtherCAT bus-OP flag, polled every 10 cycles), it could become
\texttt{true} well before the CiA402 state machine reached
\texttt{OPERATION\_ENABLED}.  The driver could declare activation success with
drives still in \texttt{SWITCH\_ON\_DISABLED}.

\begin{lstlisting}[caption={Bug~1: before and after}]
// BEFORE (wrong):
initialized_ = is_operational_;

// AFTER (correct):
initialized_ = is_operational_ && (state_ == STATE_OPERATION_ENABLED);
\end{lstlisting}

Additionally, the activation loop was simplified: the redundant
\texttt{module->initialized()} check was removed, relying solely on
\texttt{readyForCommands()} (statusword-based, every cycle) which is both
faster and more precise than the 10-cycle-polling \texttt{is\_operational\_}
flag.

\subsection{Bug~2 -- Shadowed Function Name in Barrier}
\label{sec:bug2}

\textbf{Location:} \texttt{ethercat\_driver.cpp} (unnamed namespace) vs.\
\texttt{generic\_ec\_cia402\_drive.hpp/cpp}

\textbf{Impact:} When the barrier-startup feature was added, a virtual method
\texttt{cia402PowerupRank()} was added to the \texttt{EcSlave} base class
and overridden in \texttt{EcCiA402Drive}.  An \emph{unnamed namespace} in
\texttt{ethercat\_driver.cpp} also contained a free function named
\texttt{cia402PowerupRank(int state)}.  The driver's call site called the
free function (compile-time lookup) rather than the virtual override
(runtime dispatch), meaning the rank was computed from a function that was
not aware of the current instance's state.

\textbf{Fix:} The static helper was renamed to \texttt{powerupRankOf()} and
moved to the \texttt{EcCiA402Drive} class as a \texttt{static} method.  The
unnamed namespace block in \texttt{ethercat\_driver.cpp} was deleted
entirely.  The virtual base returns $-1$ (``not a CiA402 drive'') for
non-drive slaves.

\begin{lstlisting}[caption={Bug~2: renamed static helper}]
// ec_slave.hpp (base class):
virtual int cia402PowerupRank() const { return -1; }

// generic_ec_cia402_drive.hpp:
int cia402PowerupRank() const override { return powerupRankOf(state_); }
static int powerupRankOf(DeviceState state);

// generic_ec_cia402_drive.cpp:
int EcCiA402Drive::powerupRankOf(DeviceState state) {
  switch (state) {
    case STATE_SWITCH_ON_DISABLED:  return 0;
    case STATE_READY_TO_SWITCH_ON:  return 1;
    case STATE_SWITCH_ON:           return 2;
    case STATE_OPERATION_ENABLED:   return 3;
    default:                        return -1;
  }
}
\end{lstlisting}

\subsection{Bug~3 -- Null Pointer in \texttt{transferAll()}}
\label{sec:bug3}

\textbf{Location:} \texttt{ec\_master.cpp}, \texttt{transferAll()}

\textbf{Impact:} The transfer-copy loop iterated over all registered
\texttt{EcTransfer} objects and called \texttt{memcpy(out\_ptr, in\_ptr,
size)} without checking for null pointers.  If a transfer was incompletely
configured (e.g., a module registered in the transfer table but its PDO
domain address not yet resolved), the copy caused a null-pointer dereference
and a process crash.

\begin{lstlisting}[caption={Bug~3: null-pointer guard in \texttt{transferAll()}}]
for (auto & transfer : transfers_) {
  if (transfer.in_ptr == nullptr || transfer.out_ptr == nullptr) {
    RCLCPP_ERROR(..., "transferAll: null pointer, skipping transfer");
    continue;
  }
  memcpy(transfer.out_ptr, transfer.in_ptr, transfer.size);
}
\end{lstlisting}

\subsection{Bug~4 -- Wrong Default for \texttt{phase\_stable\_cycles\_}}
\label{sec:bug4}

\textbf{Location:} \texttt{ethercat\_driver.hpp}

The default value was 20 cycles, but the design intent (matching the WC
state polling frequency and the \texttt{init\_stable\_cycles\_} default)
was 10 cycles.  With 20 the barrier phase took twice as long to pass after
all drives were ready.

\textbf{Fix:} \texttt{int phase\_stable\_cycles\_ = 20}
$\to$ \texttt{int phase\_stable\_cycles\_ = 10}.

\subsection{Bug~5 -- Missing Total-Timeout Log}
\label{sec:bug5}

\textbf{Location:} \texttt{ethercat\_driver.cpp}, start of \texttt{on\_activate()}

Without a log of the combined timeout budget (\texttt{phase\_timeout\_} +
\texttt{init\_timeout\_}), operators had no indication of the expected wait
time during startup.  Added a single \texttt{RCLCPP\_INFO} at the entry
point of \texttt{on\_activate()}.

\subsection{Bug~6 -- SDO Failure Logged at \texttt{INFO}}
\label{sec:bug6}

\textbf{Location:} \texttt{ethercat\_driver.cpp}, SDO configuration loop

When \texttt{ecrt\_master\_sdo\_download()} returned a non-zero error code,
the original code logged at \texttt{RCLCPP\_INFO} level, making SDO
misconfiguration invisible in normal operation.

\textbf{Fix:} Changed to \texttt{RCLCPP\_ERROR} and included both the
return code and the abort code in the message:

\begin{lstlisting}[caption={Bug~6: SDO error logging}]
RCLCPP_ERROR(...,
  "Failed to download SDO for module at position %s: "
  "ret=%d abort_code=0x%08X",
  ec_module_parameters_[i]["position"].c_str(),
  ret, abort_code);
\end{lstlisting}

\subsection{Bug~7 -- \texttt{std::stod} for Integer Slave Position}
\label{sec:bug7}

\textbf{Location:} \texttt{ethercat\_driver.cpp}, SDO loop

The slave position YAML parameter (always an integer like \texttt{"0"},
\texttt{"1"}, \ldots) was parsed with \texttt{std::stod} and implicitly
truncated when passed to \texttt{configSlaveSdo()} which takes
\texttt{uint16\_t}.  Floating-point parsing of ``\texttt{0}'' is fine in
practice, but the intent is clearly an unsigned integer.

\textbf{Fix:} \texttt{std::stod(...)} $\to$ \texttt{std::stoul(...)}.

\subsection{Bug~8 -- Double Push of Transfer Slaves}
\label{sec:bug8}

\textbf{Location:} \texttt{ethercat\_driver.cpp}, transfer-network
configuration

The transfer-slave index vector \texttt{ec\_transfer\_slaves\_} was built by
two loops that iterated over the same range of modules.  The first loop
pushed all indices in the range; the second loop (the surviving one) did the
same.  Every transfer slave was thus registered twice.

\textbf{Fix:} Deleted the first (duplicate) loop body entirely.

\subsection{Bug~9 -- \texttt{kMaxConsecutiveWcFailures} Not Runtime-Configurable}
\label{sec:bug9}

\textbf{Location:} \texttt{ethercat\_driver.hpp/.cpp}

The WC failure threshold was declared \texttt{static constexpr int
kMaxConsecutiveWcFailures = 5} (compile-time constant).  This meant it could
not be overridden via the hardware parameter \texttt{max\_wc\_failures} in
the URDF.  The parameter-parsing block was also missing.

\textbf{Fix:}
\begin{enumerate}[nosep]
  \item Renamed to \texttt{int kMaxConsecutiveWcFailures\_ = 5} (instance
        member with default).
  \item Added parameter parsing in \texttt{on\_configure()}:
\begin{lstlisting}
if (info_.hardware_parameters.find("max_wc_failures") !=
    info_.hardware_parameters.end()) {
  kMaxConsecutiveWcFailures_ =
    std::stoi(info_.hardware_parameters["max_wc_failures"]);
}
if (kMaxConsecutiveWcFailures_ < 1) { kMaxConsecutiveWcFailures_ = 1; }
\end{lstlisting}
\end{enumerate}

% ==============================================================================
\section{NaN Safety Analysis}
\label{sec:nan}

IEEE\,754 NaN (Not-a-Number) has a counterintuitive property: \texttt{NaN ==
x} is \emph{always} \texttt{false}, including \texttt{NaN == 0}.  C++
programmers familiar with ordinary float comparisons may write ``\texttt{if
(x == 0)} \ldots\ \texttt{else}'' logic that silently fails for NaN inputs.
In a real-time motor driver this can block safety-critical state transitions.

\subsection{NaN Sources in the Driver}

\begin{itemize}[nosep]
  \item \textbf{Startup}: All command and state interface vectors are
        initialised to \texttt{std::numeric\_limits<double>::quiet\_NaN()}.
        Until the first PDO cycle returns data, \texttt{last\_position\_} is
        NaN.
  \item \textbf{Controller inactive}: When \texttt{ros2\_control}
        deactivates a controller, it writes NaN to all command interfaces it
        owned.  This is the \emph{intended} signal that no controller owns
        the interface.
  \item \textbf{PDO decode}: \texttt{EC\_READ\_S32()} reads a signed 32-bit
        integer from the PDO domain and casts it to \texttt{double} via
        \texttt{static\_cast}, which is always finite ($\pm 2^{31}$).
        Therefore \texttt{last\_position\_} is never NaN \emph{after} the
        first PDO read.
\end{itemize}

\subsection{Existing NaN Firewall in \texttt{ec\_write()}}

The upstream channel manager already had a NaN firewall:

\begin{lstlisting}[caption={Existing NaN guard in \texttt{ec\_write()}}]
void EcPdoSingleInterfaceChannelManager::ec_write(
  uint8_t * domain_address, double value, ...)
{
  if (!std::isnan(default_value)) {
    write_function_(domain_address, last_value, mask);
  } else {
    return;  // Do nothing -- domain unchanged
  }
}
\end{lstlisting}

This prevents a NaN \texttt{default\_value} from ever being sent to the
drive.  However, it does not prevent NaN from corrupting internal state.

\subsection{NaN Bug~1 -- Fault-Reset Latch Stuck on NaN}
\label{sec:nan1}

\textbf{Symptom:} After a drive fault, commanding fault-reset (writing 1.0
to the fault-reset command interface) triggered the reset correctly.  But if
the controller was then deactivated (command interface set to NaN), the
\texttt{last\_fault\_reset\_command\_} flag remained \texttt{true} forever.
The next fault would be \emph{silently ignored}: the driver saw the
\texttt{last\_fault\_reset\_command\_} still as \texttt{true}, so no rising
edge was detected and \texttt{fault\_reset\_} was never set.

\textbf{Root cause:}
\begin{lstlisting}[caption={NaN Bug~1: original (broken) fault-reset logic}]
// BROKEN: NaN == 0 is false in IEEE 754 -- latch stays true forever
if (command_interface_ptr_->at(fault_reset_command_interface_index_) == 0) {
    last_fault_reset_command_ = false;
}
if (!last_fault_reset_command_ &&
    command_interface_ptr_->at(fault_reset_command_interface_index_) != 0) {
    last_fault_reset_command_ = true;
    fault_reset_ = true;
}
\end{lstlisting}

\textbf{Fix:}
\begin{lstlisting}[caption={NaN Bug~1: corrected fault-reset logic}]
const double fr_cmd =
  command_interface_ptr_->at(fault_reset_command_interface_index_);
// NaN == 0 is always false -> must check isnan explicitly
if (std::isnan(fr_cmd) || fr_cmd == 0) {
  last_fault_reset_command_ = false;
}
if (!last_fault_reset_command_ && !std::isnan(fr_cmd) && fr_cmd != 0) {
  last_fault_reset_command_ = true;
  fault_reset_ = true;
}
\end{lstlisting}

\subsection{NaN Bug~2 -- Position-Hold Latch Entered on NaN}
\label{sec:nan2}

\textbf{Symptom:} At startup (before the first PDO cycle),
\texttt{last\_position\_} is NaN.  The controller is inactive (command NaN),
so the position-hold branch was entered.  The condition
\texttt{!position\_hold\_initialized\_} was \texttt{true}, so
\texttt{default\_value} was set to \texttt{factor * NaN + offset = NaN}.

Although the NaN firewall in \texttt{ec\_write()} prevents the NaN from
reaching the drive, setting \texttt{default\_value = NaN} and
\texttt{position\_hold\_initialized\_ = true} caused the latch to be
considered ``initialised'' with an invalid value.  When the controller later
became active and then inactive again, the hold latch would not re-latch
from the now-valid position.

\textbf{Fix:} Guard the latch entry with \texttt{std::isfinite(last\_position\_)}:

\begin{lstlisting}[caption={NaN Bug~2: finite guard on position-hold latch}]
// BEFORE:
if (!position_hold_initialized_) {
    channel.default_value = channel.factor * last_position_ + channel.offset;
    position_hold_initialized_ = true;
}

// AFTER:
if (!position_hold_initialized_ && std::isfinite(last_position_)) {
    channel.default_value = channel.factor * last_position_ + channel.offset;
    position_hold_initialized_ = true;
}
\end{lstlisting}

\subsection{NaN Bug~3 -- Controller-Active Path Propagated NaN}
\label{sec:nan3}

\textbf{Symptom:} When the controller was active but \texttt{last\_position\_}
was still NaN (possible in the first cycle before any TPDO arrived),
\texttt{default\_value} was set to NaN unconditionally:

\begin{lstlisting}[caption={NaN Bug~3: unguarded assignment in controller-active path}]
// BEFORE (always written, even when last_position_ is NaN):
channel.default_value = channel.factor * last_position_ + channel.offset;
\end{lstlisting}

While the NaN firewall in \texttt{ec\_write()} prevents transmission, the
NaN in \texttt{default\_value} was then used to \emph{update} the latch on
the next inactive cycle, breaking the position-hold invariant.

\textbf{Fix:}
\begin{lstlisting}[caption={NaN Bug~3: finite guard on controller-active path}]
// AFTER: only update default_value when position is known-finite
if (std::isfinite(last_position_)) {
  channel.default_value = channel.factor * last_position_ + channel.offset;
}
\end{lstlisting}

\subsection{NaN Safety Summary}

The three-layer NaN defence is:
\begin{enumerate}[nosep]
  \item \textbf{Mode guard}: The entire position block is skipped if
        \texttt{mode\_of\_operation\_display\_ == MODE\_NO\_MODE} (drive not
        yet in CSP mode), preventing any write while the drive is starting up.
  \item \textbf{Finite guards (new)}: \texttt{std::isfinite(last\_position\_)}
        before every assignment to \texttt{default\_value}.
  \item \textbf{NaN firewall (existing)}: \texttt{ec\_write()} skips the write
        when \texttt{default\_value} is NaN.
\end{enumerate}

After the first PDO cycle, \texttt{EC\_READ\_S32()} always produces a finite
\texttt{double}, so \texttt{last\_position\_} is never NaN after cycle\,1.
The finite guards in layers 1 and 2 cover only the brief startup window.

% ==============================================================================
\section{Architecture of the Deterministic Barrier}
\label{sec:barrier-arch}

\subsection{CiA402 Power-Up Sequence}

CiA402 defines a mandatory sequential power-up path through the following
device states:

\begin{center}
\begin{tikzpicture}[
  state/.style={rectangle,rounded corners,draw,minimum width=2.8cm,
                minimum height=0.6cm,font=\small\ttfamily,fill=blue!7},
  arrow/.style={-Stealth,thick,blue!50}
]
  \node[state] (sod)  {SWITCH\_ON\_DISABLED};
  \node[state,right=1.4cm of sod] (rtso) {READY\_TO\_SWITCH\_ON};
  \node[state,right=1.4cm of rtso] (so)   {SWITCH\_ON};
  \node[state,right=1.4cm of so]   (oe)   {OPERATION\_ENABLED};

  \draw[arrow] (sod)  -- node[above,font=\tiny]{CW 0x06} (rtso);
  \draw[arrow] (rtso) -- node[above,font=\tiny]{CW 0x07} (so);
  \draw[arrow] (so)   -- node[above,font=\tiny]{CW 0x0F} (oe);
\end{tikzpicture}
\end{center}

The power-up ranks assigned in \texttt{powerupRankOf()} are:
\texttt{SWITCH\_ON\_DISABLED=0}, \texttt{READY\_TO\_SWITCH\_ON=1},
\texttt{SWITCH\_ON=2}, \texttt{OPERATION\_ENABLED=3}.  States outside this
path (fault, quick-stop, undefined) return $-1$.

\subsection{Barrier Mechanism}

\begin{center}
\begin{tikzpicture}[
  node distance=1.2cm,
  phase/.style={rectangle,draw,fill=orange!10,minimum width=8cm,
                minimum height=0.7cm,font=\small},
  check/.style={diamond,draw,fill=yellow!20,aspect=2.5,font=\footnotesize},
  action/.style={rectangle,draw,fill=green!10,minimum width=5cm,font=\small}
]
  \node[phase] (p0arm) {Arm barrier: all CiA402 drives set to target=SOD};
  \node[check,below=0.5cm of p0arm] (check) {WC\_COMPLETE \&\& all drives rank $\geq 0$?};
  \node[action,right=2cm of check] (stable)  {stable\_cycles++};
  \node[action,below=0.5cm of check] (wait)   {clock\_nanosleep + master->update()};
  \node[phase,below=0.5cm of wait] (release) {Disarm barrier: setStartupBarrier(false)};
  \node[phase,below=0.5cm of release] (init)  {Init gate: wait all drives OPERATION\_ENABLED};

  \draw[-Stealth] (p0arm) -- (check);
  \draw[-Stealth] (check) -- node[right,font=\tiny]{yes} (stable);
  \draw[-Stealth] (check) -- node[right,font=\tiny]{no / reset} (wait);
  \draw[-Stealth] (wait.west) -- ++(-1,0) |- (check.west);
  \draw[-Stealth] (stable.south) -- ++(0,-0.4) -| node[above right,font=\tiny]{$\geq$ stable\_cycles} (release.east);
  \draw[-Stealth] (release) -- (init);
\end{tikzpicture}
\end{center}

\textbf{Phase 0 (bus gate):} Every CiA402 drive has its barrier armed with
target state \texttt{SWITCH\_ON\_DISABLED} (rank 0).  Inside
\texttt{processData()}, when the drive reaches or passes the target rank, the
control word is held at its current value instead of advancing.  The barrier
holds until:
\begin{enumerate}[nosep]
  \item The domain WC state is \texttt{EC\_WC\_COMPLETE} (all slaves OP), \emph{and}
  \item Every CiA402 drive has power-up rank $\geq 0$ (is at or past
        \texttt{SWITCH\_ON\_DISABLED}).
\end{enumerate}

This condition must hold for \texttt{phase\_stable\_cycles\_} (10) consecutive
cycles to account for transient WC glitches on bus bring-up.

After the barrier is released, all drives transition freely through
\texttt{READY\_TO\_SWITCH\_ON $\to$ SWITCH\_ON $\to$ OPERATION\_ENABLED}
in parallel, reaching full readiness in milliseconds rather than seconds.

\subsection{Fault Safety During Barrier}

The barrier checks \texttt{module->cia402State() >= 0} before calling
\texttt{setStartupBarrier(true, target)}.  This means:
\begin{itemize}[nosep]
  \item Non-drive slaves (EK/EL Beckhoff terminals) are never armed.
  \item A drive in FAULT state has rank $-1$ and is excluded from the
        ``all at target'' check, allowing fault recovery to proceed
        independently via \texttt{transition()}.
\end{itemize}

\texttt{setStartupBarrier(false, 0)} passes 0 as the target, but
\texttt{barrier\_target\_state\_} is only read when
\texttt{barrier\_enabled\_} is \texttt{true}, so passing 0 when disarming
is safe.

% ==============================================================================
\section{Runtime Safety Supervision}
\label{sec:runtime-safety}

Once the system is activated, the \texttt{read()} callback enforces three
ongoing safety invariants:

\begin{enumerate}
  \item \textbf{Link check}: \texttt{isMasterLinkUp()} -- cable pull or
        switch power-off triggers an immediate stop.
  \item \textbf{WC check}: domain WC state must be
        \texttt{EC\_WC\_COMPLETE} for at most \texttt{kMaxConsecutiveWcFailures\_}
        (default 5) consecutive cycles.  A persistent WC fault (all slaves
        not responding) triggers a controlled stop.
  \item \textbf{Drive group failfast}: if any CiA402 drive leaves
        \texttt{OPERATION\_ENABLED} or loses a valid position for
        \texttt{runtime\_drive\_fault\_cycles\_} (default 10) consecutive
        cycles, the entire arm stops (\emph{no partial operation}).
        The faulted drive is identified by alias and position in the error
        log, with its \texttt{statusString()}.
\end{enumerate}

The rationale for group failfast (stopping all drives when one faults) is
that partial operation of a robotic arm with one joint disabled can result in
uncontrolled motion due to gravity and inertia.

% ==============================================================================
\section{Configurable Hardware Parameters}
\label{sec:params}

All safety and timing parameters are configurable from the URDF/xacro
\texttt{<hardware>} block without recompilation:

\begin{table}[H]
  \centering
  \caption{Hardware parameters and their defaults}
  \label{tab:params}
  \small
  \begin{tabularx}{\textwidth}{@{}l l X@{}}
    \toprule
    Parameter & Default & Description \\
    \midrule
    \texttt{init\_timeout} & 30\,s &
      Max time for all modules to reach OPERATION\_ENABLED \\
    \texttt{init\_stable\_cycles} & 10 &
      Cycles all modules must simultaneously be ready \\
    \texttt{phase\_timeout} & 20\,s &
      Max time for barrier phase 0 (bus gate) \\
    \texttt{phase\_stable\_cycles} & 10 &
      Cycles WC\_COMPLETE must hold before releasing barrier \\
    \texttt{startup\_mode} & \texttt{barrier} &
      \texttt{barrier} or \texttt{legacy} (per-drive auto-transition) \\
    \texttt{max\_wc\_failures} & 5 &
      Consecutive WC failures before controlled stop \\
    \texttt{runtime\_drive\_supervision} & \texttt{true} &
      Enable/disable runtime group failfast \\
    \texttt{runtime\_drive\_fault\_cycles} & 10 &
      Cycles drive fault must persist before group stop \\
    \bottomrule
  \end{tabularx}
\end{table}

\textbf{Example URDF snippet:}
\begin{lstlisting}[language=XML,caption={Hardware parameters in URDF}]
<hardware>
  <plugin>ethercat_driver/EthercatDriver</plugin>
  <param name="master_id">0</param>
  <param name="control_frequency">1000</param>
  <param name="startup_mode">barrier</param>
  <param name="phase_timeout">20</param>
  <param name="phase_stable_cycles">10</param>
  <param name="init_timeout">30</param>
  <param name="init_stable_cycles">10</param>
  <param name="max_wc_failures">5</param>
  <param name="runtime_drive_supervision">true</param>
  <param name="runtime_drive_fault_cycles">10</param>
</hardware>
\end{lstlisting}

% ==============================================================================
\section{ROS\,2 Lifecycle and Threading Model}
\label{sec:lifecycle}

\begin{center}
\begin{tikzpicture}[
  cb/.style={rectangle,rounded corners,draw,fill=blue!8,minimum width=3.8cm,
             minimum height=0.8cm,font=\small\ttfamily},
  arrow/.style={-Stealth,thick}
]
  \node[cb] (oninit)    {\textbf{on\_init()}\\param parse, alloc NaN};
  \node[cb,right=1.5cm of oninit] (oncfg)  {\textbf{on\_configure()}\\SDO, plugins};
  \node[cb,right=1.5cm of oncfg] (onact)   {\textbf{on\_activate()}\\barrier + init gate};
  \node[cb,below=1.2cm of onact] (rdwr)    {\textbf{read()/write()}\\RT loop @ 1\,kHz};
  \node[cb,below=1.2cm of oncfg] (ondact)  {\textbf{on\_deactivate()}\\master->deactivate()};
  \node[cb,below=1.2cm of oninit] (onerr)  {\textbf{on\_error()}\\stop + shutdown};

  \draw[arrow] (oninit) -- (oncfg);
  \draw[arrow] (oncfg)  -- (onact);
  \draw[arrow] (onact)  -- (rdwr);
  \draw[arrow] (rdwr.south) -- ++(0,-0.5) -| (ondact.north);
  \draw[arrow] (rdwr.south) -- ++(0,-0.5) -| (onerr.north);
  \draw[arrow] (ondact.west) -- ++(-0.5,0) |- (onerr.east);
\end{tikzpicture}
\end{center}

The real-time loop (\texttt{read()}/\texttt{write()}) runs in the
\texttt{ros2\_control} update thread at 1\,kHz.  All \texttt{on\_*()}
lifecycle callbacks hold \texttt{ec\_mutex\_} to prevent concurrent access.
The \texttt{read()} call uses \texttt{std::try\_to\_lock} so that a
lifecycle transition (which holds the mutex) does not block the RT thread
indefinitely.

% ==============================================================================
\section{Plugin Architecture}
\label{sec:plugins}

\begin{infobox}[title=Pluginlib overview]
  The driver loads EtherCAT slave implementations at runtime via
  \texttt{pluginlib}.  Every slave plugin inherits from
  \texttt{ethercat\_interface::EcSlave} and is declared via
  \texttt{PLUGINLIB\_EXPORT\_CLASS}.  The driver URDF specifies the plugin
  class name; the classloader finds the shared library at runtime.
\end{infobox}

Key virtual methods in \texttt{EcSlave}:
\begin{itemize}[nosep]
  \item \texttt{setupSlave()} -- parse YAML parameters, map PDO channels
  \item \texttt{processData(entry\_idx, domain\_address)} -- called per PDO
        entry every cycle
  \item \texttt{readyForCommands()} -- drive-level readiness (default false)
  \item \textbf{new} \texttt{hasValidPosition()} -- at least one finite
        position received (default true for non-drive slaves)
  \item \textbf{new} \texttt{statusString()} -- diagnostic string for error logs
  \item \textbf{new} \texttt{cia402State()} -- current CiA402 state as int,
        $-1$ for non-drive slaves
  \item \textbf{new} \texttt{cia402PowerupRank()} -- power-up path rank,
        $-1$ for non-drive slaves
  \item \textbf{new} \texttt{setStartupBarrier(enabled, target)} -- arm/disarm
        group barrier
\end{itemize}

The new methods have safe defaults in the base class so that non-CiA402
slaves (e.g., Beckhoff EL analogue modules) work without modification.

% ==============================================================================
\section{Summary of All Changes by File}
\label{sec:changes-by-file}

\begin{longtable}{@{}>{\ttfamily\small}p{7cm} p{8.5cm}@{}}
  \toprule
  \normalfont\textbf{File} & \normalfont\textbf{Change} \\
  \midrule
  \endfirsthead
  \multicolumn{2}{c}{\small\itshape (continued)} \\
  \toprule
  \normalfont\textbf{File} & \normalfont\textbf{Change} \\
  \midrule
  \endhead
  \bottomrule
  \endlastfoot

  ethercat\_interface/include/\\ethercat\_interface/ec\_slave.hpp &
    Added 5 new virtual methods with safe defaults:
    \texttt{readyForCommands()}, \texttt{hasValidPosition()},
    \texttt{statusString()}, \texttt{cia402State()},
    \texttt{cia402PowerupRank()}, \texttt{setStartupBarrier()} \\
  \addlinespace

  ethercat\_interface/src/ec\_master.cpp &
    Fixed per-slave random SYNC0 phase (Bug DC); added null-ptr guard in
    \texttt{transferAll()} (Bug~3); added application-time monotonic update
    in \texttt{update()}; added \texttt{deactivate()} + \texttt{release}
    for clean shutdown \\
  \addlinespace

  ethercat\_driver/include/\\ethercat\_driver/ethercat\_driver.hpp &
    Added all barrier/init/supervision member variables; changed
    \texttt{kMaxConsecutiveWcFailures} from \texttt{constexpr} to
    \texttt{int kMaxConsecutiveWcFailures\_} (Bug~9);
    \texttt{phase\_stable\_cycles\_ = 10} (Bug~4);
    declared \texttt{runBarrierStartup()} \\
  \addlinespace

  ethercat\_driver/src/ethercat\_driver.cpp &
    Added parameter parsing for all new parameters (Bugs~4,6,7,8,9);
    safety stop in \texttt{read()} (link + WC + group failfast);
    \texttt{on\_error()} lifecycle;
    \texttt{runBarrierStartup()} implementation;
    strict init gate in \texttt{on\_activate()};
    removed duplicate transfer-slave push (Bug~8);
    \texttt{stoul} for slave position (Bug~7);
    SDO errors at \texttt{RCLCPP\_ERROR} (Bug~6);
    deleted unnamed namespace with shadowed \texttt{cia402PowerupRank} (Bug~2) \\
  \addlinespace

  ethercat\_generic\_plugins/\ldots/generic\_ec\_cia402\_drive.hpp &
    Added \texttt{position\_hold\_initialized\_};
    added override declarations for all new \texttt{EcSlave} virtual methods;
    renamed static helper \texttt{cia402PowerupRank} $\to$
    \texttt{powerupRankOf} (Bug~2);
    \texttt{last\_position\_} initialised to NaN;
    added barrier member variables \\
  \addlinespace

  ethercat\_generic\_plugins/\ldots/generic\_ec\_cia402\_drive.cpp &
    Fixed \texttt{initialized\_} semantics (Bug~1);
    NaN-safe fault-reset latch (NaN Bug~1);
    finite guard on position-hold latch entry (NaN Bug~2);
    finite guard on controller-active path (NaN Bug~3);
    \texttt{powerupRankOf()} static method (Bug~2);
    added \texttt{readyForCommands()}, \texttt{hasValidPosition()},
    \texttt{statusString()} implementations \\
\end{longtable}

% ==============================================================================
\section{Commissioning and Test Plan}
\label{sec:testplan}

\subsection{Pre-Flight Checks}

\begin{enumerate}[nosep]
  \item Verify IgH EtherCAT Master is running:
        \texttt{sudo /etc/init.d/ethercat status}
  \item Confirm all 6 drives appear:
        \texttt{ethercat slaves} should show 6 entries in PREOP
  \item Confirm PREEMPT\_RT kernel: \texttt{uname -r} should show \texttt{-rt}
  \item Confirm real-time scheduling for the ROS\,2 process (CPU affinity,
        \texttt{SCHED\_FIFO})
\end{enumerate}

\subsection{Startup Test Protocol}

\begin{enumerate}
  \item Launch the driver with \texttt{startup\_mode: barrier}.
  \item Observe the log: expect\\
        \texttt{Barrier startup: phase-0 bus gate (phase\_timeout=20.0s, stable\_cycles=10)}\\
        then $<$1\,s later:\\
        \texttt{Barrier startup: all drives reached 'SWITCH\_ON\_DISABLED'}\\
        then $\sim$50\,ms later:\\
        \texttt{System Successfully started!}
  \item Total startup time should be $<$2\,s with fixed SYNC0.
  \item Compare to legacy mode (\texttt{startup\_mode: legacy}): expect
        15--20\,s with the original random-phase DC code.
\end{enumerate}

\subsection{Safety Stop Tests}

\begin{enumerate}
  \item \textbf{Cable pull test}: with the arm stationary, physically unplug
        the EtherCAT cable.  Expect the driver to log
        \texttt{SAFETY: EtherCAT link DOWN} and call \texttt{on\_error()}
        within 1\,ms.
  \item \textbf{WC failure test}: power-cycle one slave while running.
        Expect \texttt{Domain WC not COMPLETE for 5 consecutive cycles}.
  \item \textbf{Drive fault test}: send a fault-inducing command to one drive
        (overcurrent, etc.).  Expect
        \texttt{SAFETY: Drive ... left ready state ... group failfast}.
  \item \textbf{Fault-reset test}: after a drive fault, command
        \texttt{fault\_reset = 1.0} via the command interface, then
        \texttt{fault\_reset = 0.0} (or NaN).  Expect the drive to recover
        and return to OPERATION\_ENABLED.  Repeat three times to verify no
        latch-stuck regression.
\end{enumerate}

\subsection{NaN Safety Tests}

\begin{enumerate}
  \item Start the driver and immediately deactivate the trajectory controller
        (set command to NaN).  Verify no crash, no SDO errors, and that
        \texttt{last\_position\_} becomes finite within the first PDO cycle.
  \item Activate the trajectory controller, command a position, then
        deactivate.  Verify the drive holds position (no drift) and
        \texttt{position\_hold\_initialized\_} latches only once after
        \texttt{last\_position\_} is finite.
  \item Trigger a drive fault, then command \texttt{fault\_reset = 1.0} via
        the interface while the controller is active.  Set the command to NaN.
        Verify \texttt{last\_fault\_reset\_command\_} resets to \texttt{false}
        (check via a debug log or debugger watchpoint).
\end{enumerate}

% ==============================================================================
\section{Conclusion}
\label{sec:conclusion}

The evolution of \texttt{ethercat\_driver\_ros2} from the upstream
\texttt{main} branch through \texttt{saftey\_features\_clean} and into
\texttt{deterministic-barrier-startup} represents a comprehensive hardening of
the driver for use with safety-critical robotic hardware:

\begin{itemize}[nosep]
  \item \textbf{Immediate safety stops} on link-down, WC failure, and
        individual drive faults, with group-failfast semantics to prevent
        partial-arm operation.
  \item \textbf{Deterministic startup}: a single fixed SYNC0 phase eliminates
        the 15--20\,s non-deterministic DC alignment delay; the phase-0
        barrier ensures all drives advance through the CiA402 state machine
        simultaneously.
  \item \textbf{Eleven bugs fixed}, including three IEEE\,754 NaN-safety
        issues in the fault-reset and position-hold paths, a function-name
        shadow that caused the barrier rank to be computed incorrectly, and
        a null-pointer risk in the transfer subsystem.
  \item \textbf{Full runtime configurability} of all timeouts, thresholds,
        and modes via URDF hardware parameters.
\end{itemize}

The result is a driver that can be safely deployed on a robotic arm teststand
with confidence that any single-point EtherCAT failure triggers a controlled
stop within milliseconds, and that startup proceeds deterministically in under
two seconds.

% ==============================================================================
\appendix

\section{Git Commit Reference}
\label{app:commits}

\begin{longtable}{@{}>{\ttfamily\small}l >{\small}l p{9cm}@{}}
  \toprule
  \normalfont Hash & \normalfont Date & \normalfont Message \\
  \midrule
  \endfirsthead
  \multicolumn{3}{c}{\small\itshape (continued)} \\
  \toprule
  \normalfont Hash & \normalfont Date & \normalfont Message \\
  \midrule
  \endhead
  \bottomrule
  \endlastfoot
  b21391f & 2026-04-08 & safety features: immediate ERROR on link-down/WKC fault \\
  70f0595 & 2026-04-17 & ethercat\_driver: tighten safety stop and clean shutdown path \\
  32a2cc6 & 2026-04-17 & cia402: hold latched position when controller command is inactive \\
  4f1e593 & 2026-04-17 & ethercat\_interface: tune watchdog setup and use monotonic app time \\
  aedb118 & 2026-05-20 & fix(driver): raise kMaxConsecutiveWcFailures from 1 to 5 \\
  0ff09ed & 2026-06-19 & fix(dc): deterministic SYNC0 shift to eliminate per-joint init delay \\
  a8268c7 & 2026-06-19 & feat(driver): strict init readiness-gate with timeout and stability \\
  902be31 & 2026-06-24 & fix(driver): align phase-0 barrier timeout and debounce \\
  3917d66 & 2026-06-28 & WIP: fix 11 bugs in barrier startup, NaN handling, and code quality \\
\end{longtable}

\section{Key API Reference}
\label{app:api}

\begin{description}
  \item[\texttt{ecrt\_master\_state()}] Returns \texttt{ec\_master\_state\_t}
        with \texttt{link\_up}, \texttt{al\_states}, \texttt{slaves\_responding}.
  \item[\texttt{ecrt\_domain\_state()}] Returns \texttt{ec\_domain\_state\_t}
        with \texttt{working\_counter} and \texttt{wc\_state}
        (\texttt{EC\_WC\_ZERO}/\texttt{INCOMPLETE}/\texttt{COMPLETE}).
  \item[\texttt{ecrt\_slave\_config\_dc()}] Configures Distributed Clocks.
        Arguments: config, assign\_activate (0x0100=SYNC0),
        sync0\_cycle\_ns, sync0\_shift\_ns, sync1\_cycle\_ns, sync1\_shift\_ns.
  \item[\texttt{ecrt\_master\_sdo\_download()}] Blocking SDO write in
        pre-operational state.  Returns 0 on success; sets \texttt{abort\_code}.
  \item[\texttt{clock\_nanosleep(CLOCK\_MONOTONIC, TIMER\_ABSTIME, \&t, NULL)}]
        Absolute sleep to next cycle deadline; preferred over
        \texttt{nanosleep()} for jitter-free real-time loops.
\end{description}

\end{document}
